% ============================================================
%  TÓM TẮT (Tiếng Việt)
% ============================================================
\chapter*{Tóm tắt}
\addcontentsline{toc}{chapter}{Tóm tắt}

\noindent
\textbf{Từ khóa:} Truyền thông Backscatter Môi trường, chống gây nhiễu, mạng D2D, UAV relay,
học tăng cường đa tác tử, MADDPG, Imitation Learning, Graph Attention Network, Transformer.

\bigskip

Trong bối cảnh mạng truyền thông thế hệ mới, các thiết bị IoT và người dùng di động mặt đất
ngày càng đối mặt với nguy cơ bị tấn công gây nhiễu vô tuyến (jamming) có chủ đích, gây suy
giảm nghiêm trọng chất lượng liên lạc. Đề tài này nghiên cứu bài toán chống gây nhiễu cho mạng
truyền thông Device-to-Device (D2D) theo hướng tiếp cận mới: thay vì né tránh tín hiệu nhiễu,
hệ thống \emph{khai thác} năng lượng RF từ tín hiệu gây nhiễu công suất cao làm nguồn kích
hoạt cho truyền thông Ambient Backscatter (AmBC), đồng thời triển khai các UAV làm relay
động trong không gian ba chiều nhằm hỗ trợ đường truyền D2D.

Hệ thống gồm $N = 5$ cặp thiết bị nguồn--đích (SU--DU) mặt đất, $K = 2$ UAV relay di động,
một trạm relay cố định (RBS) và một thiết bị gây nhiễu mặt đất. Mỗi nguồn SU có thể chọn
một trong ba chế độ truyền: truyền thẳng D2D, relay qua RBS, hoặc relay qua UAV, đồng thời
điều chỉnh hệ số phản xạ $\alpha_i$. Các UAV tối ưu vị trí 3D theo thời gian nhằm tối đa hóa
SINR cho các SU trong phạm vi phục vụ.

Bài toán được mô hình hóa là Dec-POMDP đa tác tử gồm $N$ SU agent và $K$ UAV agent, và được
giải quyết bằng framework \textbf{IA-MADDPG mở rộng} (Imitation-Augmented Multi-Agent DDPG).
Đóng góp chính của đề tài là thay thế bộ critic MLP truyền thống bằng
\textbf{Transformer-GAT Centralized Critic}: bộ mã hóa Transformer xử lý quan sát của toàn
bộ tác tử dưới điều kiện thông tin không đầy đủ, trong khi Graph Attention Network (GAT) khai
thác cấu trúc đồ thị kết nối UAV--SU để tổng hợp thông tin topo mạng. Expert policy heuristic
được dùng để khởi động bộ nhớ replay (warm-up) và cung cấp tín hiệu bắt chước (behavior
cloning) trong giai đoạn đầu huấn luyện, nhờ đó tránh được bất ổn cold-start đặc trưng của
MARL.

Kết quả mô phỏng trên Python với 600 episode cho thấy phương pháp đề xuất vượt trội so với
năm baseline về tổng phần thưởng mạng, tỷ lệ truyền thành công (TSR), thông lượng trung
bình và hiệu quả năng lượng. Nghiên cứu ablation xác nhận đóng góp độc lập của từng thành
phần: UAV relay, imitation learning, Prioritized Experience Replay, và Transformer-GAT critic.

% ============================================================
%  TÓM TẮT TIẾNG ANH
% ============================================================
\chapter*{Abstract}
\addcontentsline{toc}{chapter}{Abstract}

\noindent
\textbf{Keywords:} truyền thông backscatter môi trường, chống gây nhiễu, mạng D2D,
UAV relay, học tăng cường sâu đa tác tử, MADDPG, imitation learning,
mạng chú ý đồ thị, Transformer.

\bigskip

Các mạng truyền thông thế hệ tiếp theo đang đối mặt ngày càng nhiều với các cuộc tấn công gây
nhiễu vô tuyến (jamming) có chủ đích, làm suy giảm nghiêm trọng chất lượng truyền thông của
các thiết bị IoT giới hạn năng lượng và người dùng di động mặt đất. Đề tài này nghiên cứu
bài toán chống gây nhiễu trong mạng truyền thông Device-to-Device (D2D) từ một góc nhìn
mới: thay vì né tránh tín hiệu jamming, hệ thống chủ động \emph{khai thác} năng lượng RF
từ tín hiệu nhiễu công suất cao làm nguồn kích hoạt cho truyền thông Ambient Backscatter
(AmBC), đồng thời triển khai các UAV làm relay động ba chiều để hỗ trợ đường truyền D2D.

Mô hình hệ thống gồm $N = 5$ cặp thiết bị nguồn--đích (SU--DU) mặt đất, $K = 2$ UAV relay
di động, một trạm relay cố định (RBS) và một thiết bị gây nhiễu mặt đất liên tục. Mỗi thiết
bị nguồn lựa chọn một trong ba chế độ truyền theo từng khe thời gian --- truyền thẳng D2D,
relay qua RBS, hoặc relay qua UAV --- đồng thời điều chỉnh hệ số phản xạ~$\alpha_i$. Các UAV
điều chỉnh vị trí 3D liên tục nhằm tối đa hóa SINR cho các thiết bị SU trong phạm vi phục vụ.

Bài toán tối ưu liên hợp được phát biểu dưới dạng Quá trình Quyết định Markov Phi Tập
trung Quan sát Một phần (Dec-POMDP) với $N$ SU agent và $K$ UAV agent, và được giải quyết
bằng framework mở rộng \textbf{IA-MADDPG} (Imitation-Augmented MADDPG). Đóng góp trung tâm
của đề tài là thay thế critic MLP truyền thống bằng \textbf{Transformer-GAT Centralized Critic}
mới: bộ mã hóa Transformer tổng hợp quan sát đa tác tử trong điều kiện thông tin không đầy
đủ, trong khi Graph Attention Network (GAT) mã hóa đồ thị topo UAV--SU để nắm bắt các phụ
thuộc không gian giữa các tác tử. Ngoài ra, expert policy được sử dụng để khởi tạo trước
bộ nhớ replay (warm-up) và cung cấp giám sát behavior cloning trong giai đoạn đầu huấn
luyện, loại bỏ bất ổn cold-start đặc trưng của MARL. Các kỹ thuật bổ sung bao gồm Prioritized
Experience Replay (PER) và làm mịn chính sách mục tiêu theo kiểu TD3.

Kết quả mô phỏng Python qua 600 episode huấn luyện cho thấy phương pháp đề xuất vượt trội
so với năm phương pháp baseline trên toàn bộ chỉ số đánh giá: phần thưởng mạng, tỷ lệ
truyền thành công (TSR), thông lượng trung bình, hiệu quả lựa chọn chế độ và hiệu quả
năng lượng. Nghiên cứu ablation định lượng đóng góp độc lập của UAV relay, imitation
learning, PER và Transformer-GAT critic.
